\documentclass[10pt]{beamer}

\newcommand{\lectnum}{L09}
\newcommand{\lecttitle}{Philosophy of AI: LLMs, Rationality and Understanding}

\usepackage{amsmath, amssymb, graphicx}
\usepackage[]{algorithm2e}
\usepackage{pdfpages}
\usepackage{bm}
\usepackage{relsize}
\usepackage[british]{babel}
\usepackage{tikz,tikzsymbols}
\usepackage{hyperref}
\usepackage{cancel}
\usepackage{caption,subcaption}
\usepackage{adjustbox}
\usepackage{overpic}
\usepackage[UKenglish]{isodate}

\usetikzlibrary{arrows,decorations.markings,decorations.pathreplacing,positioning}

\newcommand{\bs}{\boldsymbol}


\hypersetup{colorlinks,linkcolor=,urlcolor=blue}
\newenvironment{titledslide}[1]{\begin{frame}\frametitle{#1}}{\end{frame}}

\mode<presentation>{\setbeamercovered{transparent}}

\setbeamertemplate{sidebar right}{}
\setbeamertemplate{footline}{%
\hfill\usebeamertemplate***{navigation symbols}
\hspace{0.4cm}\lectnum: \insertframenumber{}/\inserttotalframenumber \hspace*{0.4cm}}

\author{James Cussens}

\title{COMS20018, Introduction to AI:\\ \vspace{5pt} \lecttitle}

\institute{School of Computer Science\\University of Bristol}
%\date{5th October 2026}

%%%%%%%%%%%%%%%%%%%%%%%%%%%%%%%%%%%%%%%%%%%%%%%%%%%%%%%%%%%%%%%%%%%%%%

\begin{document}
%%%%%%%%%%%%%%%%%%%%%%%%%%%%%%%%%%%%%%%%%%%%%%%%%%%%%%%%%%%%%%%%%%%%%%

\begin{frame}[fragile]
  \begin{tikzpicture}[remember picture, overlay]
    \node[anchor=north west, xshift=0.4cm, yshift=-0.4cm] at (current page.north west)
      {\includegraphics[width=3.4cm]{figures/UoB_RGB_24.png}};
  \end{tikzpicture}
  \titlepage
\end{frame}
%%%%%%%%%%%%%%%%%%%%%%%%%%%%%%%%%%%%%%%%%%%%%%%%%%%%%%%%%%%%%%%%%%%%%%
\begin{titledslide}{Agenda}

  \begin{itemize}
  \item We're going to look at some of the issues raised by Mike
    Woolridge in his lecture
    \href{https://www.youtube.com/watch?v=CyyL0yDhr7I}{This is not the
      AI we were promised}
  \item And some initiatives hoping to take AI further.
\end{itemize}

\end{titledslide}
%%%%%%%%%%%%%%%%%%%%%%%%%%%%%%%%%%%%%%%%%%%%%%%%%%%%%%%%%%%%%%%%%%%%%%
\begin{titledslide}{LLMs (and related systems) are game changers}

  \begin{itemize}
  \item Similarly to Woolridge's talk we'll be looking at limitations
    of LLMs etc.
  \item But we're only interested in the limitations because they are
    so impressive.
  \end{itemize}

\vspace*{2cm}
  
  \begin{quote}
    There is only one thing in the world worse than being talked
    about, and that is not being talked about.
  \end{quote} \emph{The Picture of Dorian Gray}, Oscar Wilde.
  
\end{titledslide}
%%%%%%%%%%%%%%%%%%%%%%%%%%%%%%%%%%%%%%%%%%%%%%%%%%%%%%%%%%%%%%%%%%%%%%
\begin{titledslide}{Well, what AI we're we promised?}

  \begin{itemize}
  \item It is important to understand what a revolution modern
    generative AI is!
  \item AI used to be about building `rational agents'.
  \item Contemporary AI (e.g.\ LLMs) is not like this.
  \end{itemize}
  
\end{titledslide}
%%%%%%%%%%%%%%%%%%%%%%%%%%%%%%%%%%%%%%%%%%%%%%%%%%%%%%%%%%%%%%%%%%%%%%
\begin{titledslide}{Some LLM weirdness}

  \begin{itemize}
  \item Woolridge gives many (amusing) examples or LLM irrationality.
  \item My friend Bryon Aragam also has some interesting examples from
    a NeurIPS 2024 paper. Let's have a look at it \dots
  \item As Woolridge observes: GPT-5 still ``hallucinates like
    crazy''.
  \item Looks like eliminating hallucinations properly is a tough problem.
  \end{itemize}
  
\end{titledslide}
%%%%%%%%%%%%%%%%%%%%%%%%%%%%%%%%%%%%%%%%%%%%%%%%%%%%%%%%%%%%%%%%%%%%%%
\begin{titledslide}{Sometimes the problems are more subtle}


  \begin{itemize}
  \item Gemini on ``What are Bayesian networks?'' and ``What is a
    Bayesian network?''
  \end{itemize}
  
\end{titledslide}
%%%%%%%%%%%%%%%%%%%%%%%%%%%%%%%%%%%%%%%%%%%%%%%%%%%%%%%%%%%%%%%%%%%%%%
\begin{titledslide}{LLM reasoning}

  \begin{itemize}
  \item One can, of course, show many examples of LLMs reasoning
    correctly.
  \item How do they do this?
  \item As Woolridge notes: natural language \textbf{encodes} logical
    reasoning and compositional structures. [my emphasis]
  \item But are LLMs rational minds?
  \end{itemize}
  
\end{titledslide}
%%%%%%%%%%%%%%%%%%%%%%%%%%%%%%%%%%%%%%%%%%%%%%%%%%%%%%%%%%%%%%%%%%%%%%
\begin{titledslide}{Chain-of-thought reasoning}

  \begin{itemize}
  \item Chain-of-thought reasoning is the state-of-the-art.
  \item Shifts effort to inference time.
  \item ``Looks like algorithmic intelligence'' - Woolridge.
  \end{itemize}
  
\end{titledslide}
%%%%%%%%%%%%%%%%%%%%%%%%%%%%%%%%%%%%%%%%%%%%%%%%%%%%%%%%%%%%%%%%%%%%%%
\begin{titledslide}{Robots}

  \begin{itemize}
  \item Robotic AI is the next threshold - Woolridge
  \item We still don't have robots that can do, say, household chores.
  \item So we have AI that replace some creative activities, but not
    the boring stuff - should be the other way round!
  \item Why is `physical AI' harder than `textual AI'?
  \end{itemize}
  
\end{titledslide}
%%%%%%%%%%%%%%%%%%%%%%%%%%%%%%%%%%%%%%%%%%%%%%%%%%%%%%%%%%%%%%%%%%%%%%
\begin{titledslide}{AI in the UK}

  \begin{itemize}
  \item ``UKRI launches two AI research labs to stay ahead in global
    race''
  \item \pounds 60 million in funding
  \item UK government is looking for a niche for UK AI
  \item Wants to have `AI sovereignty'
  \end{itemize}
\end{titledslide}
%%%%%%%%%%%%%%%%%%%%%%%%%%%%%%%%%%%%%%%%%%%%%%%%%%%%%%%%%%%%%%%%%%%%%%
\begin{titledslide}{SOFAIR project}
\begin{quote}
    We're very excited that UCL will be leading the new SOFAIR Lab.
    While current AI systems are impressive, many still suffer from
    basic issues such as inaccurate responses to questions. These
    systems often use similar underlying architectures, so SOFAIR will
    bring together the broader sciences and fresh ideas to create a
    new generation of open-source models. This will reduce dependency
    on the small number of model providers, boosting UK sovereignty
    and its position as a global player in AI. (David Barber, UCL)
  \end{quote}
  
\end{titledslide}
%%%%%%%%%%%%%%%%%%%%%%%%%%%%%%%%%%%%%%%%%%%%%%%%%%%%%%%%%%%%%%%%%%%%%%
\begin{titledslide}{BOLD project}

  \begin{quote}
    The UK cannot win the global AI race simply by trying to outspend
    the largest technology companies on data and compute.  BOLD is
    about a different route: discovering fundamentally new ways to
    build AI that are more efficient, more open and better aligned
    with human needs.  By focusing on new paradigms for learning,
    rather than only scaling existing methods, we aim to help secure
    the UK's sovereign capability in AI and ensure that academic
    research can shape the future of the field. (Jakob Foerster,
    University of Oxford)
  \end{quote}
  
\end{titledslide}
%%%%%%%%%%%%%%%%%%%%%%%%%%%%%%%%%%%%%%%%%%%%%%%%%%%%%%%%%%%%%%%%%%%%%%
\end{document}
